\begin{frame}{Définition}
Langage de transformation pour les documents XML
    \begin{figure}
        \centering
        \includegraphics[width=.7\linewidth]{04 - Traiter des documents XML/schema_xslt_transform.png}
    \end{figure}
\end{frame}

\begin{frame}{Utiliser XSLT dans Oxygen}
\begin{block}{Créer un fichier XSLT}
    \begin{itemize}
        \item Ouvrir le fichier \texttt{CM5\_DuBellay.xml}
        \item Ouvrir un nouveau fichier xsl (taper dans le filtre 'xsl' puis cliquer sur 'XSLT Stylesheet')
    \end{itemize}
\end{block}

\end{frame}

\begin{frame}[fragile]{En-tête xslt}
    \begin{lstlisting}[numbers=none]
        <?xml version="1.0" encoding="UTF-8"?>
        <xsl:stylesheet xmlns:xsl="http://www.w3.org/1999/XSL/Transform"
            xmlns:xs="http://www.w3.org/2001/XMLSchema"
            xmlns:math="http://www.w3.org/2005/xpath-functions/math"
            xmlns:tei="http://www.tei-c.org/ns/1.0"
            xpath-default-namespace="http://www.tei-c.org/ns/1.0"
            exclude-result-prefixes="xs math tei"
            version="3.0">
    \end{lstlisting}
\end{frame}

\begin{frame}[fragile]{Sortie xslt}
\textcolor{blue}{\texttt{<xsl:output>}}: définie le format de sortie de la transformation XSLT. Se met après l'en-tête XSLT.
    \begin{lstlisting}
    <!-- Sortie XML -->
        <xsl:output method="xml" indent="yes" encoding="UTF-8"/>
        
    <!-- Sortie HTML-->
        <xsl:output method="html" indent="yes" encoding="UTF-8"/>

    <!-- Sortie texte-->
        <xsl:output method="text" indent="yes" encoding="UTF-8"/>
    \end{lstlisting}
\end{frame}

\begin{frame}[fragile]{Utiliser XSLT dans Oxygen}
    \begin{block}{Tester des règles et voir ce qui se passe:}
\begin{lstlisting}[numbers=none]

    <!-- Premier test-->
        <xsl:template match="/">
            <xsl:apply-templates/>
        </xsl:template>
        
    <!-- Deuxième test-->
        <xsl:template match="text()"/>

    <!-- Troisième test-->
    <xsl:template match="/">
            <xsl:copy-of select="."/>
        <xsl:template/>
        
    \end{lstlisting}


\end{block}
\end{frame}

\begin{frame}[fragile]{Les règles de base de XSL}
    \begin{block}{\texttt{<xsl:template>}}
    \begin{itemize}
        \item Permet de donner les instructions xsl
        \item Utilise l'attribut \textcolor{orange}{\texttt{@match}} pour sélectionner un élément de l’arbre XML\\
        \textbf{Exemple:} \texttt{\textcolor{blue}{<xsl:template} \textcolor{orange}{match="TOTO"}\textcolor{blue}{>}}
        \item On peut ajouter du texte:\\
            \textcolor{blue}{\texttt{<xsl:template \textcolor{orange}{match="TOTO"}>}}\\
               \hspace{1cm} Voilà mon élément TOTO\\
            \textcolor{blue}{\texttt{</xsl:template>}}
    \end{itemize}
    \end{block}
    
\end{frame}

\begin{frame}[fragile]{Créer un élément}
    \begin{block}{Méthode 1:}
    \begin{lstlisting}[numbers=none]
        <xsl:template match="element">
            <p> toto </p>
        </xsl:template>
    \end{lstlisting}
    \end{block}
    
    \begin{block}{Méthode 2:}
        \begin{lstlisting}
            <xsl:template match="element">
                <xsl:element name="nouveau_element">
                    tata
                </xsl:element>
            </xsl:template>
        \end{lstlisting}
    \end{block}
    
\end{frame}

\begin{frame}[fragile]{Créer un attribut}
\begin{block}{Méthode 1:}
    \begin{lstlisting}
            <xsl:template match="element">
                <p type="attribut"> tutu </p>
            </xsl:template>
        \end{lstlisting}
\end{block}

\begin{block}{Méthode 2:}
    \begin{lstlisting}
            <xsl:template match="element">
                <xsl:element name="nouveau_element">
                    <xsl:attribute name="nouvel_attribut">
                        <xsl:text>valeur_attribut</xsl:text>
                    </xsl:attribute>
                    <xsl:text>Ici, il y avait mon élément</xsl:text>
                </xsl:element>
            </xsl:template>
        \end{lstlisting}
\end{block}
\end{frame}

\begin{frame}[fragile]{Apply-templates}
    \textcolor{blue}{\texttt{<xsl:apply-templates>}}: applique le template désigné à tous les éléments ou les enfants de l'élément sélectionné.
    \begin{block}{Exemple:}
        \begin{lstlisting}[numbers=none]
        <!-- On crée un template -->
            <xsl:template match="head">
                Title: <span style="color:#ff0000">
            <xsl:value-of select="."/></span>
            <br />
            </xsl:template>
        <!-- On applique le template -->
            <xsl:template match="text">
                <p> <xsl:apply-templates select="head"/> </p>
            </xsl:template>
        \end{lstlisting}
    \end{block}
\end{frame}